\documentclass[reqno]{amsart}
\usepackage[T1]{fontenc}
\usepackage{amsmath,amssymb,amsthm}
\usepackage{xr-hyper}
\usepackage[hidelinks]{hyperref}
\externaldocument[zr-]{20261125T001000--zeros-to-ramp-sums__workbook}
\externaldocument[zi-]{20261125T001100--inert-primes-and-estimate-b__workbook}
\externaldocument[cs-]{20261125T005500--columns-from-a-sphere__workbook}
\externaldocument[dt-]{20261126T030300--the-data__workbook}
\externaldocument[mi-]{20261126T031100--theorem-c-and-the-main-inequality__workbook}

% Chapter number; \ref{chap} gives it.
\newcounter{chap}
% Chapter numbers, generated from 20261125T012000--grh-workbook__guide.org
% by scripts/make-chapters.py.  Do not edit: change the order of the
% index and rerun the script (./preview-site.sh and CI do this).
%
% Usage in a sheet, after \newcounter{chap}:
%   % Chapter numbers, generated from 20261125T012000--grh-workbook__guide.org
% by scripts/make-chapters.py.  Do not edit: change the order of the
% index and rerun the script (./preview-site.sh and CI do this).
%
% Usage in a sheet, after \newcounter{chap}:
%   % Chapter numbers, generated from 20261125T012000--grh-workbook__guide.org
% by scripts/make-chapters.py.  Do not edit: change the order of the
% index and rerun the script (./preview-site.sh and CI do this).
%
% Usage in a sheet, after \newcounter{chap}:
%   \input{chapters}\chapnumber{<sheet ID>}
\makeatletter
\providecommand\chapters@set[2]{\@namedef{chapters@#1}{#2}}
\chapters@set{20261125T001000}{1}% zeros-to-ramp-sums
\chapters@set{20261125T001100}{2}% inert-primes-and-estimate-b
\chapters@set{20261125T005000}{3}% galois-fourier-and-averaging
\chapters@set{20261125T005500}{4}% columns-from-a-sphere
\chapters@set{20261125T000000}{5}% slope-method
\chapters@set{20261126T030100}{6}% a-twist-for-each-prime
\chapters@set{20261125T003000}{7}% the-chain-and-the-square-root
\chapters@set{20261125T010000}{8}% rank-lemma-and-zero-estimate
\chapters@set{20261125T010500}{9}% poincare-series-and-near-orthogonality
\chapters@set{20261124T223000}{10}% elliptic-curves-for-the-proof
\chapters@set{20261125T009500}{11}% fock-model-and-theta-intertwiners
\chapters@set{20261124T223100}{12}% complex-multiplication-and-inert-reduction
\chapters@set{20261125T008000}{13}% one-cm-determinant
\chapters@set{20261124T214500}{14}% supersingular-theta-divisibility
\chapters@set{20261124T214600}{15}% the-isotropic-plane
\chapters@set{20261125T013700}{16}% pairings-on-a-formal-group-of-height-two
\chapters@set{20261125T013500}{17}% hecke-nondegeneracy-on-the-tree
\chapters@set{20261126T030200}{18}% the-bound-is-attained
\chapters@set{20261125T009000}{19}% kronecker-limit-and-chowla-selberg
\chapters@set{20261125T002000}{20}% definite-theta-lifts
\chapters@set{20261125T015500}{21}% siegel-weil-in-the-smallest-case
\chapters@set{20261125T002100}{22}% definite-theta-lifts-and-norms
\chapters@set{20261125T016500}{23}% hecke-operators-and-theta-lifts
\chapters@set{20261125T016700}{24}% divisibility-along-the-chain
\chapters@set{20261125T004000}{25}% denominators-and-sha
\chapters@set{20261125T004100}{26}% tate-twists-and-sha
\chapters@set{20261126T020430}{27}% the-minus-part-by-gauss-sums
\chapters@set{20261125T014100}{28}% congruence-modules
\chapters@set{20261125T014500}{29}% eisenstein-congruences-and-selmer-classes
\chapters@set{20261125T014700}{30}% pseudocharacters-and-the-length-bound
\chapters@set{20261125T011000}{31}% galois-representations-black-box
\chapters@set{20261126T030300}{32}% the-data
\chapters@set{20261126T020210}{33}% bessel-poincare-series-and-their-gram-operators
\chapters@set{20261126T030400}{34}% the-seesaw-identity
\chapters@set{20261126T030500}{35}% divisibility-of-the-composite-lift
\chapters@set{20261126T030600}{36}% theorem-a-short-forms
\chapters@set{20261126T030700}{37}% the-determinant-of-a-chain
\chapters@set{20261126T030800}{38}% the-cost-of-one-lift
\chapters@set{20261126T030900}{39}% denominators-of-projections-to-lifts
\chapters@set{20261126T031000}{40}% twists-with-few-congruences
\chapters@set{20261126T031100}{41}% theorem-c-and-the-main-inequality
\providecommand\chapnumber[1]{%
  \@ifundefined{chapters@#1}%
    {\PackageError{chapters}{Sheet #1 is not listed in the index}%
      {List it in 20261125T012000--grh-workbook__guide.org\MessageBreak
       and run scripts/make-chapters.py.}}%
    {\edef\thechap{\@nameuse{chapters@#1}}}}
\makeatother
\chapnumber{<sheet ID>}
\makeatletter
\providecommand\chapters@set[2]{\@namedef{chapters@#1}{#2}}
\chapters@set{20261125T001000}{1}% zeros-to-ramp-sums
\chapters@set{20261125T001100}{2}% inert-primes-and-estimate-b
\chapters@set{20261125T005000}{3}% galois-fourier-and-averaging
\chapters@set{20261125T005500}{4}% columns-from-a-sphere
\chapters@set{20261125T000000}{5}% slope-method
\chapters@set{20261126T030100}{6}% a-twist-for-each-prime
\chapters@set{20261125T003000}{7}% the-chain-and-the-square-root
\chapters@set{20261125T010000}{8}% rank-lemma-and-zero-estimate
\chapters@set{20261125T010500}{9}% poincare-series-and-near-orthogonality
\chapters@set{20261124T223000}{10}% elliptic-curves-for-the-proof
\chapters@set{20261125T009500}{11}% fock-model-and-theta-intertwiners
\chapters@set{20261124T223100}{12}% complex-multiplication-and-inert-reduction
\chapters@set{20261125T008000}{13}% one-cm-determinant
\chapters@set{20261124T214500}{14}% supersingular-theta-divisibility
\chapters@set{20261124T214600}{15}% the-isotropic-plane
\chapters@set{20261125T013700}{16}% pairings-on-a-formal-group-of-height-two
\chapters@set{20261125T013500}{17}% hecke-nondegeneracy-on-the-tree
\chapters@set{20261126T030200}{18}% the-bound-is-attained
\chapters@set{20261125T009000}{19}% kronecker-limit-and-chowla-selberg
\chapters@set{20261125T002000}{20}% definite-theta-lifts
\chapters@set{20261125T015500}{21}% siegel-weil-in-the-smallest-case
\chapters@set{20261125T002100}{22}% definite-theta-lifts-and-norms
\chapters@set{20261125T016500}{23}% hecke-operators-and-theta-lifts
\chapters@set{20261125T016700}{24}% divisibility-along-the-chain
\chapters@set{20261125T004000}{25}% denominators-and-sha
\chapters@set{20261125T004100}{26}% tate-twists-and-sha
\chapters@set{20261126T020430}{27}% the-minus-part-by-gauss-sums
\chapters@set{20261125T014100}{28}% congruence-modules
\chapters@set{20261125T014500}{29}% eisenstein-congruences-and-selmer-classes
\chapters@set{20261125T014700}{30}% pseudocharacters-and-the-length-bound
\chapters@set{20261125T011000}{31}% galois-representations-black-box
\chapters@set{20261126T030300}{32}% the-data
\chapters@set{20261126T020210}{33}% bessel-poincare-series-and-their-gram-operators
\chapters@set{20261126T030400}{34}% the-seesaw-identity
\chapters@set{20261126T030500}{35}% divisibility-of-the-composite-lift
\chapters@set{20261126T030600}{36}% theorem-a-short-forms
\chapters@set{20261126T030700}{37}% the-determinant-of-a-chain
\chapters@set{20261126T030800}{38}% the-cost-of-one-lift
\chapters@set{20261126T030900}{39}% denominators-of-projections-to-lifts
\chapters@set{20261126T031000}{40}% twists-with-few-congruences
\chapters@set{20261126T031100}{41}% theorem-c-and-the-main-inequality
\providecommand\chapnumber[1]{%
  \@ifundefined{chapters@#1}%
    {\PackageError{chapters}{Sheet #1 is not listed in the index}%
      {List it in 20261125T012000--grh-workbook__guide.org\MessageBreak
       and run scripts/make-chapters.py.}}%
    {\edef\thechap{\@nameuse{chapters@#1}}}}
\makeatother
\chapnumber{<sheet ID>}
\makeatletter
\providecommand\chapters@set[2]{\@namedef{chapters@#1}{#2}}
\chapters@set{20261125T001000}{1}% zeros-to-ramp-sums
\chapters@set{20261125T001100}{2}% inert-primes-and-estimate-b
\chapters@set{20261125T005000}{3}% galois-fourier-and-averaging
\chapters@set{20261125T005500}{4}% columns-from-a-sphere
\chapters@set{20261125T000000}{5}% slope-method
\chapters@set{20261126T030100}{6}% a-twist-for-each-prime
\chapters@set{20261125T003000}{7}% the-chain-and-the-square-root
\chapters@set{20261125T010000}{8}% rank-lemma-and-zero-estimate
\chapters@set{20261125T010500}{9}% poincare-series-and-near-orthogonality
\chapters@set{20261124T223000}{10}% elliptic-curves-for-the-proof
\chapters@set{20261125T009500}{11}% fock-model-and-theta-intertwiners
\chapters@set{20261124T223100}{12}% complex-multiplication-and-inert-reduction
\chapters@set{20261125T008000}{13}% one-cm-determinant
\chapters@set{20261124T214500}{14}% supersingular-theta-divisibility
\chapters@set{20261124T214600}{15}% the-isotropic-plane
\chapters@set{20261125T013700}{16}% pairings-on-a-formal-group-of-height-two
\chapters@set{20261125T013500}{17}% hecke-nondegeneracy-on-the-tree
\chapters@set{20261126T030200}{18}% the-bound-is-attained
\chapters@set{20261125T009000}{19}% kronecker-limit-and-chowla-selberg
\chapters@set{20261125T002000}{20}% definite-theta-lifts
\chapters@set{20261125T015500}{21}% siegel-weil-in-the-smallest-case
\chapters@set{20261125T002100}{22}% definite-theta-lifts-and-norms
\chapters@set{20261125T016500}{23}% hecke-operators-and-theta-lifts
\chapters@set{20261125T016700}{24}% divisibility-along-the-chain
\chapters@set{20261125T004000}{25}% denominators-and-sha
\chapters@set{20261125T004100}{26}% tate-twists-and-sha
\chapters@set{20261126T020430}{27}% the-minus-part-by-gauss-sums
\chapters@set{20261125T014100}{28}% congruence-modules
\chapters@set{20261125T014500}{29}% eisenstein-congruences-and-selmer-classes
\chapters@set{20261125T014700}{30}% pseudocharacters-and-the-length-bound
\chapters@set{20261125T011000}{31}% galois-representations-black-box
\chapters@set{20261126T030300}{32}% the-data
\chapters@set{20261126T020210}{33}% bessel-poincare-series-and-their-gram-operators
\chapters@set{20261126T030400}{34}% the-seesaw-identity
\chapters@set{20261126T030500}{35}% divisibility-of-the-composite-lift
\chapters@set{20261126T030600}{36}% theorem-a-short-forms
\chapters@set{20261126T030700}{37}% the-determinant-of-a-chain
\chapters@set{20261126T030800}{38}% the-cost-of-one-lift
\chapters@set{20261126T030900}{39}% denominators-of-projections-to-lifts
\chapters@set{20261126T031000}{40}% twists-with-few-congruences
\chapters@set{20261126T031100}{41}% theorem-c-and-the-main-inequality
\providecommand\chapnumber[1]{%
  \@ifundefined{chapters@#1}%
    {\PackageError{chapters}{Sheet #1 is not listed in the index}%
      {List it in 20261125T012000--grh-workbook__guide.org\MessageBreak
       and run scripts/make-chapters.py.}}%
    {\edef\thechap{\@nameuse{chapters@#1}}}}
\makeatother
\chapnumber{20261125T005000}
\renewcommand{\thesection}{\thechap.\arabic{section}}

\theoremstyle{definition}
\newtheorem{problem}{Problem}[section]
\theoremstyle{plain}
\newtheorem{theorem}{Theorem}[section]

\newcommand{\Z}{\mathbb{Z}}
\newcommand{\Q}{\mathbb{Q}}
\newcommand{\C}{\mathbb{C}}
\newcommand{\R}{\mathbb{R}}
\newcommand{\Gal}{\operatorname{Gal}}
\newcommand{\rem}{\mathrm{rem}}
\newcommand{\smain}{s_{\mathrm{main}}}

\setlength{\emergencystretch}{2em}

\title{The places, the weights and the remainder}
\author{}
\date{November 26, 2026}

\begin{document}
\maketitle
\refstepcounter{chap}\label{chap}

\paragraph{Prerequisites.}
From Chapter~\ref{zr-chap}, \emph{From zeros to a one-sided bound}: the weight \(\lambda\), the constant \(c_{d,q}\) and the remainder \(R_{d,q}\) of equation~\eqref{zr-eq:R}, recalled below, and Problem~\ref{zr-pr:R-limit}.  From Chapter~\ref{zi-chap}, \emph{Inert primes and the choice of \(d\)}: Section~\ref{zi-sec:inert}, on the values of \(1 - \chi_K\).  Cyclotomic fields and their Galois groups \cite[Ch.~2]{Wa}; infinite places of a number field as conjugation orbits of embeddings into \(\C\) \cite[Ch.~II, \S8]{Ne}; quadratic reciprocity \cite[Ch.~I, \S3]{SeA}.

\medskip
In the proof, the main inequality (Theorem~\ref{mi-thm:main}) bounds from above a centered sum over the infinite places of class \(1\) modulo \(q_0\) in the CM field obtained by adjoining a root of unity to \(K = \Q(\sqrt{-d})\).  Chapter~\ref{zr-chap} needs an upper bound for the remainder of a weighted sum over the prime powers \(n \equiv 1 \pmod{q_0}\) up to \(x\).  How does the sum over places give such a bound?  An exact finite-sum identity gives the answer.  In the summand indexed by \(n\), the weight at the translated place sets the cutoff.  For an admissible datum the centered total over places of class \(1\) modulo \(q_0\) equals a positive multiple of \(R_{q_0}(x)\) plus two terms at \(x_1\) (Theorem~\ref{thm:remainder}).

Throughout, \(d \equiv 3 \pmod 4\) is a prime, \(K := \Q(\sqrt{-d})\), and \(\chi_K(n) := \bigl(\frac nd\bigr)\) is the Legendre symbol (the \(\chi_d\) of Chapter~\ref{zr-chap}), so \(\chi_K(d) = 0\).  For an integer \(q \ge 1\) prime to \(d\) we write, as in Chapter~\ref{zr-chap},
\[
  \lambda_q(n) := \begin{cases} \Lambda(n)\bigl(1 - \chi_K(n)\bigr), & n \equiv 1 \pmod q, \\ 0, & \text{otherwise}, \end{cases}
  \qquad
  R_q(z) := \sum_{1 \le n \le z} \lambda_q(n) \Bigl(\frac1n - \frac1z\Bigr) - \frac{\log z}{\varphi(q)} - c_q
\]
for real \(z > 0\), where \(c_q := c_{d,q}\) and \(R_q := R_{d,q}\) are the constant and the remainder of equation~\eqref{zr-eq:R}, and \(\lambda_q\) is the weight \(\lambda\) there for the modulus \(q\); we drop \(d\), which is fixed.  The value of \(c_q\) plays no role in the identities of this chapter.  We write \(u_+ := \max(u, 0)\) for real \(u\).

\section{The host field and its places}\label{sec:places}

Let \(q_0 \ge 3\) and let \(q_1\) be a multiple of \(q_0\) such that every prime factor of \(q_1\) divides \(q_0\), and \(d \nmid q_1\).  Let \(E := K(\zeta)\) with \(\zeta\) a primitive \(q_1\)-th root of unity, regarded as a subfield of \(\C\).  For \(g \in \Gal(E/K)\) let \(e(g) \in (\Z/q_1\Z)^\times\) be the class with \(g(\zeta) = \zeta^{e(g)}\).

\begin{problem}\label{pr:galois}
The map \(e \colon \Gal(E/K) \to (\Z/q_1\Z)^\times\) is an isomorphism, and \([E : \Q] = 2\varphi(q_1)\).
\end{problem}

Let \(\Phi\) be the set of infinite places of \(E\) and \(f := \lvert \Phi \rvert\).  For \(g \in \Gal(E/K)\) let \(\langle g \rangle \in \Phi\) be the place of the embedding \(g \colon E \to \C\).

\begin{problem}\label{pr:cm-type}
Every infinite place of \(E\) is complex, and \(g \mapsto \langle g \rangle\) is a bijection \(\Gal(E/K) \to \Phi\).  Hence \(f = \varphi(q_1)\).
\end{problem}

We make \(\Phi\) a group through \(g \mapsto \langle g \rangle\), and identify it with \((\Z/q_1\Z)^\times\) through \(e\): the place \(\sigma \in \Phi\) has \emph{label} \(v \in (\Z/q_1\Z)^\times\) if \(\sigma = \langle g \rangle\) with \(e(g) = v\).  The place with label \(1\) is the identity.  For \(\nu \in \Z\) let \([\nu] \in \Phi\) be the place with label \(\nu \bmod q_1\) if \(\gcd(\nu, q_1) = 1\), and the identity place otherwise.  Let \(\mathrm{cl} \colon \Phi \to (\Z/q_0\Z)^\times\) be the label reduced modulo \(q_0\), and let
\[
  \Phi_a := \mathrm{cl}^{-1}(1) ,
\]
the \emph{\(a\)-places}.  Fix a set \(\Phi_b \subseteq \mathrm{cl}^{-1}(-1)\) of \emph{\(b\)-places}.  We write \(\#\) for the number of elements.

\begin{problem}\label{pr:fibres}
\(\#\Phi_a = q_1/q_0\) and \(f = \#\Phi_a\, \varphi(q_0)\).
\end{problem}

The place attached to an integer \(n \ge 1\) and an \(a\)-place \(\tau\) is \([-n]\tau\): for \(n\) prime to \(q_1\), its label is \(-n\) times the label of \(\tau\).

\begin{problem}\label{pr:outside}
For \(\tau \in \Phi_a\) and every integer \(n \ge 1\), \(\mathrm{cl}\bigl([-n]\tau\bigr) = -1\) if and only if \(n \equiv 1 \pmod{q_0}\).
\end{problem}

\begin{problem}\label{pr:count}
For every integer \(n \ge 1\), the number of \(\tau \in \Phi_a\) with \([-n]\tau \in \Phi_b\) is \(\#\Phi_b\) if \(n \equiv 1 \pmod{q_0}\), and \(0\) otherwise.
\end{problem}

When \(\gcd(n, q_1) > 1\), our convention gives \([-n]\tau = \tau\).  Under the hypotheses of the next problem such \(n\) have coefficient \(\lambda_1(n) = 0\) in the prime sums below, so the convention does not affect them.

\begin{problem}\label{pr:coef}
(i) If \(16 \mid q_0\), \(q_0 \mid 2(d + 1)\) and \(q_0 \nmid d + 1\), then \(d \equiv 7 \pmod 8\) and \(\chi_K(p) = 1\) for every prime \(p \mid q_0\).
(ii) In that case \(\lambda_1(n) = 0\) for every \(n \ge 1\) with \(\gcd(n, q_1) > 1\).
\end{problem}

The \emph{toy example} is \(q_0 = 4\), \(q_1 = 16\), \(d = 7\), so that \(f = 8\), \(\Phi_a\) consists of the places with labels \(1, 5, 9, 13\), and \(\Phi_b\) consists of the place with label \(15\), so \(\#\Phi_a = 4\) and \(\#\Phi_b = 1\).  It does not satisfy the hypotheses of Problem~\ref{pr:coef}(i), but \(\chi_K(2) = (2/7) = 1\), and the conclusion of Problem~\ref{pr:coef}(ii) holds.

\begin{problem}\label{pr:table}
In the toy example, the primes \(p \le 61\) with \(\lambda_1(p) \ne 0\) are \(7\) and the primes inert in \(K\), and the label of \([-p]\) is as follows.
\[
  \begin{array}{c|cccccccc}
    \text{label} & 1 & 3 & 5 & 7 & 9 & 11 & 13 & 15 \\ \hline
    p & 31,\, 47 & 13,\, 61 & 59 & 41 & 7 & 5 & 3,\, 19 & 17
  \end{array}
\]
The other \(n \le 61\) with \(\lambda_1(n) \ne 0\) are \(27\) and \(49\), with \(\lambda_1(27) = 2\log 3\) and \(\lambda_1(49) = \log 7\).
\end{problem}

\section{Two weights and the prime sum}\label{sec:thick}

Let \(x < x_1\) and \(b\) be positive integers with \(x_1 \mid b\), and let \(a := xb\).  Define the \emph{weight} of a place \(\sigma\) by
\[
  \kappa(\sigma) := \begin{cases} a/x = b, & \sigma \in \Phi_b, \\ a/x_1, & \sigma \notin \Phi_b, \end{cases}
\]
an integer; the \(b\)-places carry the larger weight.  For real \(y\), define the \emph{prime sum} at a place \(\sigma\) by
\[
  s_{\mathrm{fin}}(y, \sigma) := \sum_{n = 1}^{a} \lambda_1(n) \Bigl( \frac yn - \kappa\bigl([-n]\sigma\bigr) \Bigr)_+ .
\]
For real \(z > 0\) and \(n \ge 1\) let \(H_z(n) := (a/n - a/z)_+\).

\begin{problem}\label{pr:ramp}
\(H_z(n) = a(1/n - 1/z)\) for \(n < z\), and \(H_z(n) = 0\) for \(n \ge z\).  Hence for \(1 \le n \le a\) the \(n\)-th term of \(s_{\mathrm{fin}}(a, \sigma)\) is \(\lambda_1(n) H_x(n)\) if \([-n]\sigma \in \Phi_b\), and \(\lambda_1(n) H_{x_1}(n)\) otherwise.
\end{problem}

\begin{problem}\label{pr:thin}
Let \(\sigma \in \Phi\).
(i) For \(0 \le y \le a\) and \(x_1 \le n \le a\), the \(n\)-th term of \(s_{\mathrm{fin}}(y, \sigma)\) is \(0\).
(ii) For \(x \le n < x_1\), the \(n\)-th term of \(s_{\mathrm{fin}}(a, \sigma)\) is \(0\) if \([-n]\sigma \in \Phi_b\), and \(\lambda_1(n)\, a (1/n - 1/x_1)\) otherwise, which is positive unless \(\lambda_1(n) = 0\).
\end{problem}

The first step towards the total over the \(a\)-places is a ramp sum in a progression.

\begin{problem}\label{pr:ramp-sum}
For integers \(q \ge 1\) prime to \(d\), \(A_0 \ge 1\), and real \(0 < z \le A_0\),
\[
  \sum_{n = 1}^{A_0} \lambda_q(n) \Bigl( \frac{A_0}{n} - \frac{A_0}{z} \Bigr)_+ = A_0 \Bigl( R_q(z) + \frac{\log z}{\varphi(q)} + c_q \Bigr) .
\]
For \(d = 7\), \(q = 4\), \(A_0 = 60\), \(z = 20\), both sides equal \(18 \log 5 + \frac{42}{13} \log 13 + \frac{18}{17} \log 17 = 40.2565\ldots\).
\end{problem}

In the average over the \(a\)-places, the place \([-n]\tau\) selects the residue class of \(n\).

\begin{problem}\label{pr:select}
Let \(q_0 = 10\), \(q_1 = 20\), \(x = 2\), \(x_1 = 3\), \(b = 6\) and \(\Phi_b = \{[9]\}\), so that \(a = 12\), \(\Phi_a = \{[1], [11]\}\), and \(\kappa\) is \(6\) at \([9]\) and \(4\) elsewhere.  For \(1 \le n \le 12\) and \(\tau \in \Phi_a\), the term \((60/n - \kappa([-n]\tau))_+\) differs from \((60/n - 4)_+\) only for \(n \in \{1, 11\}\), and then for one \(\tau\) only: at \(n = 1\) it is \(54\) instead of \(56\) (for \(\tau = [11]\)), at \(n = 11\) it is \(0\) instead of \(16/11\) (for \(\tau = [1]\)).  At \(n = 3\) both terms are \(16\).
\end{problem}

\begin{problem}\label{pr:thick-avg}
For every integer \(n \ge 1\),
\[
  \sum_{\tau \in \Phi_a} \lambda_1(n) \Bigl( \frac an - \kappa\bigl([-n]\tau\bigr) \Bigr)_+
  = \#\Phi_a\, \lambda_1(n) H_{x_1}(n) + \#\Phi_b\, \lambda_{q_0}(n) \bigl( H_x(n) - H_{x_1}(n) \bigr) .
\]
\end{problem}

\begin{problem}\label{pr:sum-A}
\[
  \sum_{\tau \in \Phi_a} s_{\mathrm{fin}}(a, \tau)
  = a \Bigl[ \#\Phi_b \bigl( R_{q_0}(x) - R_{q_0}(x_1) \bigr) + \#\Phi_a R_1(x_1) \Bigr]
  + a \Bigl[ \#\Phi_a (\log x_1 + c_1) + \#\Phi_b\, \frac{\log (x/x_1)}{\varphi(q_0)} \Bigr] .
\]
\end{problem}

\begin{problem}\label{pr:toy-sum}
In the toy example with \(x = 20\), \(x_1 = 60\), \(b = 960\), so that \(a = 19200\) and the weights are \(960\) and \(320\), the four prime sums \(s_{\mathrm{fin}}(a, \tau)\) at the places \(\tau\) with labels \(1, 5, 9, 13\) rounded to the nearest hundredth, are \(45108.01\), \(45591.18\), \(47772.92\), \(46814.23\).  Their total and the right side of Problem~\ref{pr:sum-A} both round to \(185286.35\).
\end{problem}

The expected value uses the mean of the logarithms of the weights over all places.  We define the \emph{expected value}
\[
  \smain := a \Bigl( \log a - \frac1f \sum_{\sigma \in \Phi} \log \kappa(\sigma) + c_1 \Bigr) .
\]

\begin{problem}\label{pr:log-mean}
\(\sum_{\sigma \in \Phi} \log \kappa(\sigma) = f \log (a/x_1) - \#\Phi_b \log (x/x_1)\).  In the toy example both sides equal \(47.24518\ldots\).
\end{problem}

\begin{problem}\label{pr:main}
\(\#\Phi_a\, \smain = a \bigl[ \#\Phi_a (\log x_1 + c_1) + \#\Phi_b \log(x/x_1)/\varphi(q_0) \bigr]\).
\end{problem}

Define the \emph{total remainder} by \(\rem := \sum_{\tau \in \Phi_a} \bigl( s_{\mathrm{fin}}(a, \tau) - \smain \bigr)\), so that the mean of \(s_{\mathrm{fin}}(a, \tau)\) over \(\Phi_a\) is \(\smain + \rem/\#\Phi_a\).  Subtracting the identity of Problem~\ref{pr:main} from that of Problem~\ref{pr:sum-A} gives
\begin{equation}\label{eq:rem}
  \rem = a \Bigl[ \#\Phi_b \bigl( R_{q_0}(x) - R_{q_0}(x_1) \bigr) + \#\Phi_a R_1(x_1) \Bigr] .
\end{equation}
In the toy example, \(\rem + 76800\, c_1 = -118612.63\ldots\).

In the proof the \emph{parameters} \((q_0, d)\) are a positive multiple \(q_0\) of \(18480\) and a prime \(d\) with \(q_0 \mid 2(d + 1)\) and \(q_0 \nmid d + 1\) (Definition~\ref{dt-def:parameters} of Chapter~\ref{dt-chap}).  An \emph{admissible datum} for them is a datum satisfying all the conditions of Definition~\ref{dt-def:data} there.  We record the components used here.  The scales \(x\), \(x_1\), \(b\) are positive integers.  The number \(m\) of theta lifts in each chain is a power of \(2\) with \(\sqrt x \le m \le 2\sqrt x\), and the tower exponent is \(s = 22 \log_2 m\) (logarithm to base \(2\)).  With \(e_s := \lfloor \log_3 s \rfloor\) (logarithm to base \(3\)), let \(q_1 := q_0\,3^{e_s+5}2^s\); the field \(E\), its embedding into \(\C\), the primitive root \(\zeta\) and the labels of the places are as in \S\ref{sec:places}.  We have \(q_1 x_1 \mid b\).  The datum supplies \(\Phi_b \subseteq \mathrm{cl}^{-1}(-1)\) with \(q_1/q_0 \le \#\Phi_b \log b\).  Since \(x_1 \mid b\), the prescribed weight \(\lfloor xb/x_1\rfloor\) outside \(\Phi_b\) equals \(a/x_1\).

\begin{theorem}\label{thm:remainder}
For all parameters \((q_0, d)\) and every admissible datum for them, \eqref{eq:rem} holds, where \(a = xb\) and \(R_q\) and \(c_1\) are those of Chapter~\ref{zr-chap} for the prime \(d\).
\end{theorem}

For these data \(q_0\) and \(q_1\) have the same prime factors, \(d \nmid q_1\), and \(d \equiv 7 \pmod 8\).  The scale conditions give \(0 < \log x\) and \(10^6 \log x \le \log x_1\), so \(x < x_1\).

\section{The remainder at the first scale}\label{sec:two-scales}

Chapter~\ref{zr-chap} needs an upper bound for \(R_{q_0}(x)\).  When \(\Phi_b\) is nonempty, equation~\eqref{eq:rem} isolates \(R_{q_0}(x)\) at the price of two terms at \(x_1\).

\begin{problem}\label{pr:solve}
If \(\#\Phi_b > 0\), then \(R_{q_0}(x) = \dfrac{\rem}{a\,\#\Phi_b} + R_{q_0}(x_1) - \dfrac{\#\Phi_a}{\#\Phi_b} R_1(x_1)\).
\end{problem}

For \(\#\Phi_b>0\), the coefficient of \(R_1(x_1)\) has magnitude \(\#\Phi_a/\#\Phi_b\), the reciprocal of the proportion of the \(-1\) fibre occupied by \(\Phi_b\); this is why \(\Phi_b\) must be dense.

\begin{problem}\label{pr:dense}
If \(b > 1\) and \(\#\Phi_a \le \#\Phi_b \log b\), then \(\#\Phi_b > 0\) and
\[
  R_{q_0}(x) \le \frac{\rem}{a\,\#\Phi_b} + \lvert R_{q_0}(x_1) \rvert + \log b \, \lvert R_1(x_1) \rvert .
\]
In the toy example \(\#\Phi_a/\#\Phi_b = 4 \le \log 960 = 6.8669\ldots\).
\end{problem}

Write \(\theta := \#\Phi_b/\#\Phi_a\) in this section.  The main inequality \(m \cdot \rem \le fa(\log b)^5\) bounds the mean excess over the \(a\)-places in units of \(a\), with a saving of \(1/m\).  Under the density hypothesis of Problem~\ref{pr:dense}, isolating \(R_{q_0}(x)\) costs a factor \(1/\theta \le \log b\), which gives the sixth logarithm in the bound at \(x\).

\begin{problem}\label{pr:mean-main}
For a positive integer \(m\), the main inequality implies
\[
  R_1(x_1) + \theta \bigl(R_{q_0}(x) - R_{q_0}(x_1)\bigr)
  \le \frac{\varphi(q_0)(\log b)^5}{m}.
\]
\end{problem}

\begin{problem}\label{pr:with-main}
Suppose \(b > 1\), \(\#\Phi_a \le \#\Phi_b \log b\), and \(m\) is a positive integer.  If \(m \cdot \rem \le f a (\log b)^5\), then
\[
  R_{q_0}(x) \le \frac{\varphi(q_0) (\log b)^6}{m} + \lvert R_{q_0}(x_1) \rvert + \log b \, \lvert R_1(x_1) \rvert .
\]
\end{problem}

Chapter~\ref{mi-chap} proves the main inequality for every admissible datum (Theorem~\ref{mi-thm:main}) and applies Problem~\ref{pr:with-main} in Problem~\ref{mi-pr:Rx}, before passing to the asymptotic estimate of Theorem~\ref{mi-thm:one-sided}; Chapter~\ref{cs-chap} constructs a set \(\Phi_b\) dense enough for Problem~\ref{pr:dense}.

\paragraph{In the Lean formalization.}
In the Lean proof \cite{GL}, the parameters and the admissible data are \texttt{New.Parameters} and \texttt{New.Data}; Theorem~\ref{thm:remainder} is \texttt{New.Data.totalRemainder\_eq}, in \texttt{New/Reduction/PrimeSums.lean}; and the conclusion of Problem~\ref{pr:with-main} for every admissible datum, without further hypotheses, is \texttt{New.Data.primeRemainder\_le}, in \texttt{New/Reduction/OneSidedBound.lean}.

\section*{Hints}

\subsection*{\S\ref{sec:places}}
\begin{description}
\item[\ref{pr:galois}] \(d\) ramifies in \(K\) and not in \(\Q(\zeta)\), so \(K \cap \Q(\zeta) = \Q\) \cite[Ch.~2]{Wa}.
\item[\ref{pr:cm-type}] Complex conjugation restricts to the nontrivial automorphism of \(K\).
\item[\ref{pr:fibres}] \(\mathrm{cl}\) is a surjective homomorphism, and \(\varphi(q_1)/\varphi(q_0) = q_1/q_0\) when \(q_0\) and \(q_1\) have the same prime factors.
\item[\ref{pr:outside}] \(\mathrm{cl}\) is a homomorphism; if \(\gcd(n, q_1) > 1\), then \(n \not\equiv 1 \pmod{q_0}\).
\item[\ref{pr:count}] Multiplication by \([-n]\) is a bijection of \(\Phi\); Problem~\ref{pr:outside}.
\item[\ref{pr:coef}] (i) \(d + 1 \equiv q_0/2 \pmod{q_0}\); for each odd prime \(p \mid q_0\), \(d \equiv -1 \pmod p\), and reciprocity.  (ii) Section~\ref{zi-sec:inert}.
\item[\ref{pr:table}] \(p\) is inert in \(\Q(\sqrt{-7})\) exactly when \(p \equiv 3, 5, 6 \pmod 7\).
\end{description}

\subsection*{\S\ref{sec:thick}}
\begin{description}
\item[\ref{pr:thin}] (i) Compare \(y/n\) and \(\kappa([-n]\sigma)\) with \(a/x_1\).  (ii) Use the two possible weights at \([-n]\sigma\).
\item[\ref{pr:ramp-sum}] Problem~\ref{pr:ramp} with \(A_0\) in place of \(a\), and the definition of \(R_q\).
\item[\ref{pr:select}] \([-n]\tau = [9]\) exactly when \(n\tau \equiv 11 \pmod{20}\).
\item[\ref{pr:thick-avg}] Write \((a/n - \kappa(\sigma))_+\) as \(H_{x_1}(n)\) plus a correction supported on \(\Phi_b\); Problem~\ref{pr:count}, and \(\lambda_{q_0}(n) = \lambda_1(n)\) for \(n \equiv 1 \pmod{q_0}\).
\item[\ref{pr:sum-A}] Exchange the sums over \(\tau\) and \(n\); Problem~\ref{pr:ramp-sum} three times, with \((q, z) = (1, x_1), (q_0, x), (q_0, x_1)\).
\item[\ref{pr:toy-sum}] Only the \(n < 60\) with \(\lambda_1(n) \ne 0\) contribute (Problems~\ref{pr:table} and~\ref{pr:thin}).
\item[\ref{pr:main}] Problems~\ref{pr:fibres} and~\ref{pr:log-mean}.
\end{description}

\subsection*{\S\ref{sec:two-scales}}
\begin{description}
\item[\ref{pr:solve}] Equation~\eqref{eq:rem}.
\item[\ref{pr:dense}] \(\#\Phi_a > 0\).
\item[\ref{pr:mean-main}] Divide by \(ma\,\#\Phi_a\) and use \eqref{eq:rem} and Problem~\ref{pr:fibres}.
\item[\ref{pr:with-main}] Use Problem~\ref{pr:mean-main} and the density bound on \(1/\theta\).
\end{description}

\subsection*{Theorems}
\begin{description}
\item[\ref{thm:remainder}] Compare Problems~\ref{pr:sum-A} and~\ref{pr:main}.
\end{description}

\begin{thebibliography}{SeA}

\bibitem[Ne]{Ne}
J.~Neukirch, \emph{Algebraic Number Theory}, Grundlehren Math. Wiss. \textbf{322}, Springer, 1999.

\bibitem[SeA]{SeA}
J.-P.~Serre, \emph{A Course in Arithmetic}, Grad. Texts in Math. \textbf{7}, Springer, 1973.

\bibitem[GL]{GL} L.~Sernau, \emph{GRH for Dirichlet \(L\)-functions, in Lean~4}, a shorter rewriting of the Lean proof, text by Claude, \url{https://github.com/lsernau/grh-lean}, commit \texttt{72fead1}, readable edition \texttt{grh/}.

\bibitem[Wa]{Wa}
L.~C.~Washington, \emph{Introduction to Cyclotomic Fields}, 2nd ed., Grad. Texts in Math. \textbf{83}, Springer, 1997.

\end{thebibliography}

\end{document}
